\documentclass[11pt]{article}

% --- Series style ---
\usepackage{series}

% --- Math and layout ---
\usepackage{amsmath,amssymb,amsthm}
\usepackage{mathtools}
\usepackage{geometry}
\usepackage{hyperref}

\geometry{margin=1in}


% --- Metadata ---
\title{A Tiered Roadmap for Calculations in\\
Modal Triplet Theory (MTT)}
\author{Peter Nero}
\date{January 2026}

\begin{document}
\maketitle

\begin{abstract}
We present a tiered methodology for extracting predictions from the Modal
Triplet Theory (MTT) without requiring a full internal metric solve.
The program is organized into four operational tiers with explicit
contracts: inputs allowed, tools employed, and outputs guaranteed.
Tier~1 delivers exact, topology--only results independent of geometry.
Tier~2 adds geometry--light relations and inequalities.
Tier~3 treats MTT as a \emph{superset} theory and backfits a minimal set
of latent constants from overlaps among gauge, gravitational, and
auxiliary sectors using one (or two) empirical inputs.
Tier~4 introduces a string--lift, mapping MTT bundles to heterotic or
type~IIB/F--theory compactifications so that internal data reduce to
intersection numbers, Chern classes, periods, and standard threshold
formulae.
We give explicit algorithms, uncertainty accounting, and a consolidated
test matrix that separates what is proved, what is executed numerically,
and what remains open.
This paper defines the computational architecture of the MTT program and
serves as the methodological reference for subsequent results.
\end{abstract}

\tableofcontents
\newpage

% ============================================================
\section{Introduction and Scope}

Modal Triplet Theory (MTT) realizes the Standard Model and gravity as
emergent sectors of a higher--dimensional modal geometry with three
interlocking bundle factors.
The theory has matured to the point where a wide range of exact,
bounded, and calibratable predictions can be obtained \emph{without}
solving the internal metric problem directly.

The purpose of this paper is \emph{not} to present new phenomenological
results.
Instead, it provides a \emph{roadmap} that organizes the calculational
program into operational tiers with clearly defined contracts.
Each tier specifies:
\begin{itemize}
  \item what inputs are permitted,
  \item what mathematical tools may be used,
  \item and what class of outputs is guaranteed.
\end{itemize}

This separation serves three goals:
\begin{enumerate}
  \item to deliver exact or bounded results as early as possible,
  \item to make explicit which conclusions are rigid and which are
        execution--dependent,
  \item to prevent confusion between necessity (structure) and existence
        (realization).
\end{enumerate}

Later papers in the MTT series instantiate particular tiers.
This paper defines the framework against which those results should be
read and evaluated.
\paragraph{Numerical provenance.}
All explicit numerical results that were previously distributed across standalone
numerics notes have been fully subsumed into the corresponding Tier~3 and Tier~4
execution papers. The present Roadmap defines methodology, tier contracts, and
validation criteria only; it does not duplicate executed results. Concretely:
Tier~3 numerics (crossing scales, $\zeta$-ratios, $K$, minimum-norm thresholds,
and $\alpha_s(M_Z)$ cross-predictions) appear in \emph{Superset Determinations in
Modal Triplet Theory (MTT)}, and Tier~4 numerics (K\"ahler moduli, thresholds,
axions, Yukawas, CKM/PMNS, and Higgs boundary data) appear in
\emph{Execution of Modal Triplet Theory I} and \emph{Execution of Modal Triplet
Theory II}. Earlier numerics notes may therefore be retired without loss of content.

% ============================================================
\subsection{Latent parameters and conventions}

We work on a spacetime of the form
\[
M_{10} = Y_4 \times X_6 ,
\]
where $Y_4$ is a smooth, oriented, spin four--manifold and $X_6$ is the
internal modal space.

The central latent quantities used throughout the program are:
\begin{itemize}
  \item a common gauge scale
  \[
  K := \frac{4\pi\,\mathrm{Vol}(X_6)}{g_{10}^2} ,
  \]
  \item dimensionless harmonic normalization weights
  \[
  \alpha_r^{-1}(\Lambda) \equiv \frac{4\pi}{g_r^2(\Lambda)} = K\,\zeta_r,
  \qquad r \in \{1,2,3\},
  \]
  \item and the gravitational overlap
  \[
  G_{\mathrm{eff}}^{-1} = G_{10}^{-1}\,\mathrm{Vol}(X_6),
  \qquad G_{\mathrm{eff}} \equiv G_N .
  \]
\end{itemize}

The $\zeta_r$ are defined only up to an overall normalization absorbed
into $K$; only ratios $\zeta_a/\zeta_b$ are intrinsic.
Hypercharge is always taken in GUT normalization,
$g' = \sqrt{3/5}\,g_1$.

% ============================================================
\section{Tier Taxonomy and Contracts}

We partition the MTT calculational program into four operational tiers.
Each tier has a strict contract specifying allowed inputs, tools, and
outputs.

\begin{definition}[Tier contracts]
A \emph{tier} is defined by:
\begin{enumerate}
  \item the class of admissible inputs,
  \item the mathematical tools permitted,
  \item the type of outputs guaranteed,
  \item and the accuracy class of those outputs.
\end{enumerate}
Later tiers may use results from earlier tiers but not vice versa.
\end{definition}

% ------------------------------------------------------------
\subsection{Tier~0: Units and conventions}

Tier~0 fixes units ($c=\hbar=k_B=1$), normalization conventions, and
notation.
It produces no physical predictions.

% ------------------------------------------------------------
\subsection{Tier~1: Topology--only predictions (exact)}

\paragraph{Inputs.}
Topology of $Y_4$, spin structure, Chern classes of line bundles, flux
balance conditions, and representation assignments.

\paragraph{Tools.}
Index theory, cohomology, group theory, anomaly bookkeeping, and
selection rules.

\paragraph{Outputs.}
Exact, discrete statements independent of any internal metric:
\begin{itemize}
  \item family number by Dirac index,
  \item exact Standard Model hypercharges from difference charges,
  \item cancellation of all local gauge and gravitational anomalies,
  \item absence of the SU(2) Witten global anomaly for three families,
  \item topological criterion for Dirac vs.\ Majorana masses,
  \item strong--CP relaxation via a central PQ--like symmetry with integer
        domain--wall number,
  \item forbiddance of many baryon/lepton violating operators,
  \item equality of low--energy wave speeds $c_{\mathrm{grav}}=c_{\mathrm{em}}$,
  \item holonomy determinant/phase sum rules.
\end{itemize}

\paragraph{Accuracy class.}
Exact (topological).

% ------------------------------------------------------------
\subsection{Tier~2: Geometry--light relations (bounded)}

\paragraph{Inputs.}
Tier~1 structure plus mild symmetry or positivity assumptions (e.g.\
modal democracy, spectral gaps).

\paragraph{Tools.}
Algebraic identities, inequalities, sign analysis, and first--principles
bounds.

\paragraph{Outputs.}
Relations and bounds that do not require harmonic norms:
\begin{itemize}
  \item high--scale electroweak mixing
        $\sin^2\theta_W = 3/8$ under modal symmetry,
  \item holonomy phase/determinant sum constraints,
  \item qualitative renormalization--group sign structure,
  \item curvature--mass drift identities and FRW bounds,
  \item PPN bound $\gamma = 1 + O(\Delta_{\mathrm{curv}})$.
\end{itemize}

\paragraph{Accuracy class.}
Exact identities or rigorous inequalities.

% ============================================================
\subsection{Tier~3: Calibratable constants via superset overlaps}

Tier~3 treats MTT as a \emph{superset theory}.
Instead of computing internal harmonic norms directly, one introduces a
small set of latent scalars and determines them algebraically from
overlaps among gauge, gravitational, and auxiliary sectors using a
single matching.

\paragraph{Inputs.}
\begin{itemize}
  \item Low--energy gauge couplings evolved to a matching scale
        $\Lambda_{\mathrm{MTT}}$ via a specified RGE scheme (SM or MSSM,
        one-- or two--loop).
  \item Newton's constant $G_N$.
  \item One additional non--geometric closure (modal democracy, PQ
        prior, or PPN bound).
\end{itemize}

\paragraph{Latent variables.}
\[
\mathcal{L} =
\left\{
\frac{\zeta_2}{\zeta_1},
\frac{\zeta_3}{\zeta_1},
K,
\frac{\mathrm{Vol}(X_6)}{g_{10}^2},
\mathrm{Vol}(X_6) G_{10}^{-1}
\right\}.
\]

\paragraph{Tools.}
Renormalization--group evolution, algebraic ratios, overlap identities,
and linear error propagation.

\paragraph{Outputs.}
\begin{itemize}
  \item Geometry--independent extraction of $\zeta$--ratios from data.
  \item Calibration of $K$ from a single coupling.
  \item Determination of two independent combinations of
        $(\mathrm{Vol}, g_{10}, G_{10})$.
  \item Cross--predictions at $\Lambda_{\mathrm{MTT}}$
        (e.g.\ $\sin^2\theta_W(\Lambda)$, $g_3^2/g_2^2$).
  \item Round--trip consistency checks under RGE flow.
\end{itemize}

\paragraph{Accuracy class.}
Percent--level, controlled by RGE scheme and threshold systematics.

% ------------------------------------------------------------
\subsubsection{Algorithm 1: Overlap--Backfit}

\begin{enumerate}
  \item Choose a matching scale $\Lambda_{\mathrm{MTT}}$ and RGE scheme.
  \item Run $(\alpha_1,\alpha_2,\alpha_3)$ from $M_Z$ to
        $\Lambda_{\mathrm{MTT}}$.
  \item Extract ratios
        \[
        \frac{\zeta_2}{\zeta_1}
        = \frac{\alpha_2^{-1}(\Lambda)}{\alpha_1^{-1}(\Lambda)},
        \qquad
        \frac{\zeta_3}{\zeta_1}
        = \frac{\alpha_3^{-1}(\Lambda)}{\alpha_1^{-1}(\Lambda)}.
        \]
  \item Calibrate $K = \alpha_r^{-1}(\Lambda)/\zeta_r$ using any $r$.
  \item Use $G_N$ to fix $\mathrm{Vol}(X_6) G_{10}^{-1}$.
  \item Apply one closure to solve $(\mathrm{Vol}, g_{10}, G_{10})$.
  \item Run back to $M_Z$ and verify agreement within uncertainties.
\end{enumerate}

% ============================================================
\subsection{Tier~4: String--lift geometry for quantitative predictions}

Tier~4 introduces a concrete geometric realization.
The defining principle is that internal harmonic integrals are replaced
by standard algebraic--geometric data.

\paragraph{Inputs.}
\begin{itemize}
  \item Tier~3 latents $(\zeta$--ratios, $K)$.
  \item Choice of embedding: heterotic line--bundle sum or
        type~IIB/F--theory D7--brane stacks.
\end{itemize}

\paragraph{Tools.}
Intersection rings $\kappa_{ijk}$, Chern classes, line--bundle slopes,
Picard--Fuchs periods, flux and tadpole constraints, one--loop threshold
formulae.

\paragraph{Outputs.}
\begin{itemize}
  \item K\"ahler moduli ratios and divisor volumes.
  \item Gauge kinetic matrix and axion decay constants.
  \item High--scale threshold profiles consistent with Tier~3 targets.
  \item (Optionally) Yukawa textures, CKM/PMNS, Higgs boundary data.
\end{itemize}

\paragraph{Accuracy class.}
Model--dependent execution with explicit uncertainty bands.

% ------------------------------------------------------------
\subsubsection{Algorithm 2: String--Lift Backfit}

\begin{enumerate}
  \item Choose a pilot Calabi--Yau (or toroidal) compactification with
        three distinguished divisors or line bundles satisfying
        $c_1$--balance.
  \item Compute intersection numbers $\kappa_{ijk}$ and topological
        invariants.
  \item Solve for K\"ahler moduli ratios from the Tier~3 $\zeta$--ratios.
  \item Calibrate $K$ using one gauge coupling and $G_N$.
  \item Include one--loop thresholds and check consistency with Tier~3
        minimal or minimum--norm profiles.
  \item Compute selected Yukawa entries or axion couplings as desired.
\end{enumerate}
% ============================================================
\section{Toolbox and Validation Workflows}

This section summarizes the computational tools associated with each
tier and the validation checks that ensure internal consistency across
the program.

% ------------------------------------------------------------
\subsection{Tier--wise tool summary}

\paragraph{Tier~1 tools.}
\begin{itemize}
  \item Index theorems (Atiyah--Singer).
  \item Chern class algebra and flux balance.
  \item Group--theoretic anomaly sums.
  \item Line--bundle triviality and selection rules.
\end{itemize}

\paragraph{Tier~2 tools.}
\begin{itemize}
  \item Algebraic identities (phase/determinant sums).
  \item Renormalization--group sign analysis.
  \item Curvature--mass inequalities and PPN bounds.
\end{itemize}

\paragraph{Tier~3 tools.}
\begin{itemize}
  \item One-- and two--loop RGE solvers (SM/MSSM, piecewise).
  \item Algebraic ratio extraction and calibration.
  \item Least--squares or minimum--norm threshold diagnostics.
  \item Linear error propagation and small MCMC fits (optional).
\end{itemize}

\paragraph{Tier~4 tools.}
\begin{itemize}
  \item Algebraic geometry: intersection rings $\kappa_{ijk}$.
  \item Line--bundle cohomology and slope stability.
  \item Picard--Fuchs solvers for periods.
  \item Flux, tadpole, and anomaly consistency checks.
  \item One--loop threshold formulae.
\end{itemize}

% ------------------------------------------------------------
\subsection{Validation and cross--checks}

At each tier, internal consistency is enforced by explicit checks:
\noindent\textbf{Validation and cross--checks.}
\begin{itemize}
  \item Internal algebra:
        indices, anomaly cancellation, and phase sum rules.
  \item RGE round--trip:
        parameters extracted at $\Lambda_{\mathrm{MTT}}$ are run back
        to $M_Z$ and compared with data
        (see Tier~3 Superset Determinations, Sec.~19).
  \item Normalization consistency:
        the scale $K$ agrees when extracted from different gauge factors.
  \item String--lift closure:
        slopes, tadpoles, and massless $U(1)$ conditions are satisfied,
        and threshold profiles remain controlled.
\end{itemize}


% ============================================================
\section{TOE Test Matrix}

The following matrix summarizes the status of major tests organized by
tier. ``Proved'' denotes a rigorous derivation from MTT axioms; 
``Executed'' denotes a completed numerical or algebraic realization;
``Planned'' denotes a well--specified pipeline not executed here.

\medskip

\noindent\textbf{Tier~1 (Topology--only).}
\begin{itemize}
  \item Family number by index \hfill Proved
  \item Exact SM hypercharges \hfill Proved
  \item Local and global anomaly cancellation \hfill Proved
  \item Dirac vs.\ Majorana criterion \hfill Proved
  \item Strong--CP relaxation, integer $N_{\mathrm{DW}}$ \hfill Proved
  \item Operator forbiddance \hfill Proved
  \item $c_{\mathrm{grav}} = c_{\mathrm{em}}$ \hfill Proved
\end{itemize}

\noindent\textbf{Tier~2 (Geometry--light relations).}
\begin{itemize}
  \item High--scale electroweak mixing identity
        $\sin^2\theta_W = 3/8$ under modal democracy, with controlled
        sensitivity to deviations.
  \item Holonomy determinant and phase sum rules implied by canonical
        trivialization.
  \item Qualitative renormalization--group sign structure
        (recalled from Tier~1 for completeness).
  \item Curvature--mass drift identities and FRW bounds.
  \item Post--Newtonian inequality target
        $\gamma = 1 + O(\Delta_{\mathrm{curv}})$, understood as a
        conditional Tier--2 bound.
\end{itemize}

\noindent\textbf{Tier~3 (Superset calibration).}
\begin{itemize}
  \item Geometry--independent extraction of modal weight ratios
        $\zeta_2/\zeta_1$ and $\zeta_3/\zeta_1$ from executed RGE running.
  \item Executed determination of crossing scales
        $\Lambda_{12}$ and $\Lambda_{23}$ (SM two--loop canonical).
  \item Numerical calibration of the common scale $K$ (Planck units),
        consistent across gauge factors.
  \item Determination of the two invariant combinations
        $\mathrm{Vol}(X_6)/g_{10}^2$ and
        $\mathrm{Vol}(X_6)\,G_{10}^{-1}$.
  \item Executed minimum--norm threshold diagnostic with a single bulk
        coefficient $\delta = -25.2 \pm 0.5$.
  \item Cross--prediction of $\alpha_s(M_Z)$ without additional tuning.
  \item Round--trip RGE consistency checks within stated uncertainties.
\end{itemize}


\noindent\textbf{Tier~4 (String--lift execution).}
\begin{itemize}
  \item Existence of a consistent Calabi--Yau corner realizing the
        Tier~3 invariants (proved at the level of explicit construction).
  \item Executed gauge and K\"ahler sectors:
        explicit $t_i$, $\tau_i$, and $\mathrm{Vol}(X_6)$.
  \item Executed axion sector:
        explicit decay constants $f_a$ with ratios matching
        $\zeta$--ratios.
  \item Executed reproduction of Tier~3 minimum--norm thresholds from
        geometry, including the bulk coefficient $\delta=-25.2$ and
        small exceptional--cycle corrections.
  \item Executed flavor and Higgs sectors:
        explicit Yukawa matrices, CKM matrix and Jarlskog invariant,
        PMNS matrix and neutrino masses, and Higgs quartic boundary data.
  \item Explicit uncertainty accounting and EFT control checks.
\end{itemize}


% ============================================================
% ============================================================
\section{Roadmap and Milestones}

With the tiered structure now fully instantiated, the Modal Triplet
Theory (MTT) program admits a clear and stable execution roadmap.
Each milestone below corresponds to a completed or well-defined tier,
with explicit success criteria and canonical paper references.

\noindent\textbf{M0: Methodological foundation (completed).}
\begin{itemize}
  \item Definition of the tier taxonomy, contracts, and validation
        criteria.
  \item Separation of topology--only results from geometry--light,
        calibratable, and execution--level results.
  \item Consolidation of all numerical provenance into Tier~3 and
        Tier~4 execution papers.
\end{itemize}

\noindent\textbf{M1: Tier~1 exact constraints (completed).}
\begin{itemize}
  \item Topology--only derivation of Standard Model hypercharges,
        family number, and anomaly cancellation.
  \item Dirac versus Majorana mass criterion and SU(2) global anomaly
        test.
  \item Canonical trivialization and holonomy phase sum rules.
\end{itemize}

\noindent\textbf{M2: Tier~2 geometry--light relations (completed).}
\begin{itemize}
  \item High--scale electroweak mixing identity
        $\sin^2\theta_W=3/8$ under modal democracy.
  \item Curvature--mass drift identities and FRW bounds.
  \item Conditional post--Newtonian inequality
        $\gamma = 1 + O(\Delta_{\mathrm{curv}})$.
\end{itemize}

\noindent\textbf{M3: Tier~3 superset determinations (completed).}
\begin{itemize}
  \item Geometry--independent extraction of $\zeta$--ratios and the
        common scale $K$ from executed RGE running.
  \item Executed determination of crossing scales
        $\Lambda_{12}$ and $\Lambda_{23}$.
  \item Minimum--norm threshold diagnostic with a single bulk
        coefficient $\delta=-25.2\pm0.5$.
  \item Cross--prediction of $\alpha_s(M_Z)$ and round--trip consistency
        checks.
\end{itemize}

\noindent\textbf{M4: Tier~4 execution I --- gauge, axion, and thresholds
(completed).}
\begin{itemize}
  \item Explicit Calabi--Yau corner realizing all Tier~3 invariants.
  \item Executed K\"ahler moduli, divisor volumes, and internal volume.
  \item Explicit axion decay constants and geometric reproduction of the
        Tier~3 threshold structure.
\end{itemize}

\noindent\textbf{M5: Tier~4 execution II --- flavor and Higgs sectors
(completed).}
\begin{itemize}
  \item Executed quark and lepton Yukawa matrices.
  \item Explicit CKM and PMNS matrices with CP violation.
  \item Neutrino masses and splittings from a seesaw realization.
  \item Executed Higgs quartic boundary condition consistent with
        $m_h\simeq125$~GeV.
\end{itemize}

\noindent\textbf{M6: Future directions (planned).}
\begin{itemize}
  \item Proton decay operators and lifetime estimates.
  \item Moduli stabilization and vacuum selection.
  \item Cosmological dynamics and early--universe applications.
\end{itemize}


% ============================================================
\section{Conclusions and Outlook}

We have presented a fully tiered and execution--complete framework for
calculations in Modal Triplet Theory (MTT).
All results formerly distributed across standalone numerics notes have
been consolidated into the main text of the Tier~3 and Tier~4 execution
papers, leaving the present Roadmap as a purely methodological guide.

The tier structure cleanly separates:
\begin{itemize}
  \item exact, topology--only constraints (Tier~1),
  \item geometry--light identities and bounds (Tier~2),
  \item geometry--free calibration of latent parameters (Tier~3),
  \item and explicit string--lift realizations (Tier~4).
\end{itemize}

This separation clarifies which results are rigid consequences of the
theory and which depend on particular realizations.
In particular, Tier~3 establishes MTT as a superset framework in which
all gauge and gravitational invariants are fixed algebraically from data
prior to any geometric input.
Tier~4 then demonstrates that these invariants can be realized
simultaneously in explicit string compactifications, including realistic
flavor and Higgs sectors, without retuning upstream inputs.

The resulting program is coherent, non--regressive, and falsifiable:
failure at any tier would isolate the source of inconsistency.
Conversely, successful execution through Tier~4 shows that topology,
effective field theory, and string geometry can be integrated into a
single calculational pipeline.

Future work will focus on extending the Tier~4 execution to global
consistency conditions, vacuum selection, and cosmology.
These developments will refine, but not alter, the structural results
established here.

The tiered Roadmap provided in this paper is intended as the canonical
reference for navigating the MTT corpus and for assessing future
extensions of the theory.
\begin{thebibliography}{99}

\bibitem{MTT-Admissibility}
P.~Nero,
\emph{Modal Triplet Theory: Admissibility, Encodings, and the Structure of Physical Description},
Zenodo preprint, January 2026.
\url{https://doi.org/10.5281/zenodo.18255621}

\bibitem{MTT-Foundation}
P.~Nero,
\emph{Modal Triplet Theory: Foundation},
Zenodo preprint, September 2025.
\url{https://doi.org/10.5281/zenodo.16949762}

\bibitem{MTT-FixedPoints}
P.~Nero,
\emph{Fixed Points I--VI: Complete Coherence Spine},
Zenodo preprints, August 2025.
\url{https://doi.org/10.5281/zenodo.16948748}

\bibitem{MTT-ProjectionAdmissibility}
P.~Nero,
\emph{The Projection--Admissibility Principle: Structural Constraints on Effective Physical Description},
Zenodo preprint, January 2026.
\url{https://doi.org/10.5281/zenodo.18255838}

\bibitem{MTT-Closure}
P.~Nero,
\emph{Closure and Inevitability in Modal Triplet Theory},
Zenodo preprint, January 2026.
\url{https://doi.org/10.5281/zenodo.18255510}

\bibitem{MTT-CoherenceCapacity}
P.~Nero,
\emph{Coherence Capacity as the Fundamental Resource of Effective Physics},
Zenodo preprint, January 2026.
\url{https://doi.org/10.5281/zenodo.18255905}

\bibitem{MTT-CoherenceDynamics}
P.~Nero,
\emph{Dynamics of Coherence Capacity: Transport, Concentration, and Exhaustion},
Zenodo preprint, January 2026.
\url{https://doi.org/10.5281/zenodo.18256048}

\bibitem{MTT-QM}
P.~Nero,
\emph{Modal Triplet Theory: From MTT to Quantum Mechanics},
Zenodo preprint, September 2025.
\url{https://doi.org/10.5281/zenodo.17074246}

\bibitem{MTT-QFT}
P.~Nero,
\emph{From MTT to Quantum Field Theory},
Zenodo preprint, 2025.
\url{https://doi.org/10.5281/zenodo.17068816}

\bibitem{MTT-GR}
P.~Nero,
\emph{Modal Triplet Theory: From MTT to General Relativity},
Zenodo preprint, October 2025.
\url{https://doi.org/10.5281/zenodo.16950597}

\bibitem{MTT-UVQG}
P.~Nero,
\emph{Modal Triplet Theory: From MTT to a UV-Finite, Unitary Quantum Gravity},
Zenodo preprint, 2025.
\url{https://doi.org/10.5281/zenodo.17077671}

\bibitem{MTT-Measurement}
P.~Nero,
\emph{Measurement as Disturbance and Stabilization in Modal Triplet Theory},
Zenodo preprint, 2025.
\url{https://doi.org/10.5281/zenodo.17177404}

\bibitem{MTT-Irreversibility}
P.~Nero,
\emph{Projection, Probability, and Irreversibility: Shadow Bridges Between Measurement, Black Holes, and Cosmology in Modal Triplet Theory},
Zenodo preprint, January 2026.
\url{https://doi.org/10.5281/zenodo.18256408}

\bibitem{MTT-Bell-Beables}
P.~Nero,
\emph{Modal Fixed Points, Bell’s Beables, and the Limits of Factorization},
Zenodo preprint, 2025.
\url{https://doi.org/10.5281/zenodo.17076300}

\bibitem{MTT-Bell-Temporal}
P.~Nero,
\emph{Temporal Bell Inequalities and Global Consistency in Modal Triplet Theory},
Zenodo preprint, August 2025.
\url{https://doi.org/10.5281/zenodo.18208884}

\bibitem{MTT-Stochastic}
P.~Nero,
\emph{From Modal Triplet Theory to Indivisible Stochastic Processes: A First-Principles, Fully Rigorous Derivation},
Zenodo preprint, January 2026.
\url{https://doi.org/10.5281/zenodo.18254862}

\end{thebibliography}

\end{document}
